\subsection{"追--躲"循环的识别与中断}

在意愿不一致的伴侣中，最常见的互动模式是"追--躲"（pursuer--distancer）循环：一方更频繁地发起，另一方更频繁地回避。若这一模式固化为"我总在要、你总在躲"，双方都会积累委屈——追求方感到持续被拒，回避方感到持续被逼。中断这一循环的关键，是把讨论焦点从"谁该改变"转向"我们的互动是如何形成的"，并借助前述协商与排期机制，把发起与回应从个人之间的对抗，转化为关系之内的共同安排。
